\section{Structured causal models: restrictions, computation and the target}
\label{exec23:structured-scope}

A universal bank observes accounts within branches, repeated decisions within
customers and exposures connected through markets or payment networks. The
observations are neither an arbitrary collection of independent rows nor
automatically interchangeable. Two recent primary contributions clarify
different parts of this problem. Geometric causal models use specified
symmetries to organize dependence and identification. A hierarchical causal
pipeline translates distributional interventions into symbolic expressions
and fitted numerical factors. Both provide useful ingredients for a proposed
product, but neither permits an agent to infer valid causal meaning merely
from a table schema or a successful program execution.
\citep{exec23_geometric_models,exec23_hierarchical_pipeline}

\subsection{A symmetry is a scientific restriction}
\label{exec23:structured-symmetry}

In Weinstein and Blei's geometric causal model, an endogenous variable
$x^v$ is a structured collection $(x^v_\omega)_{\omega\in\Omega}$,
where $\Omega$ indexes locations, times, nodes or other positions. A
group $G$ acts on these collections. A causal mechanism $f^v$ is equivariant
when applying a transformation before evaluating the mechanism is equivalent
to transforming its output:
\begin{equation}
 f^v(\phi(x^{\mathrm{pa}(v)}),\phi(\epsilon^v))
 =\phi\!\left(f^v(x^{\mathrm{pa}(v)},\epsilon^v)\right),
 \qquad \phi\in G.
 \label{exec23:structured-equivariance}
\end{equation}
Here $\mathrm{pa}(v)$ denotes the parents of variable $v$ in a specified
directed acyclic graph and $\epsilon^v$ is exogenous noise. The model uses
equivariant mechanisms to create dependence among endogenous positions from
the stipulated noise structure. A hard intervention replaces the selected
mechanism with a supplied value. Permutation, spatial-shift and array
symmetries express different models; they are not different software
optimizations of one unrestricted model.
\citep{exec23_geometric_models}

For a banking application, translation of a branch-level model across
locations would assert that the relevant mechanism survives that
transformation after the declared context is represented. A policy change,
regulatory boundary or market break may make that assertion implausible.
The proposal should therefore ask which transformations preserve the
decision mechanism, what variables encode exceptions, and what evidence
could falsify the proposed symmetry. These are proposed elicitation questions.
The paper assumes the causal graph and symmetry group are known and leaves
their discovery to future work. An LLM may help formulate candidate
restrictions and collect supporting documents; this paper supplies no
experiment establishing that an LLM can validate them.
\citep{exec23_geometric_models}

\subsection{Identification and consistency retain their assumptions}
\label{exec23:structured-identification}

The paper separates learning an observational law from one structured
realization, identifying an intervention from that law, and estimating it
from finitely observed coordinates. Its mixing-action condition requires
countably infinite index and transformation sets, with transformations that
eventually move any finite set away from itself. Under the stated model,
ergodicity then makes learning from an expanding observation domain possible.
The do-calculus soundness argument uses the graph's Markov decomposition.
For groups consisting of index transformations, the completeness argument
embeds conventional unrestricted causal models into the geometric class.
Stronger symmetries can add restrictions: the orthogonal-group instrumental
variable example invokes linear-Gaussian structure, a nonconstant relevant
instrument and finite interference. These conclusions have different
premises and should not be compressed into a claim that symmetry identifies
any latent graph.
\citep{exec23_geometric_models}

The finite-data consistency result is likewise conditional. It assumes
finite interference, local information under the transformations, an
amenable group with a tempered F{\o}lner sequence, finite state spaces and
positivity of the required finite observational configurations. Informally,
averages over the permitted transformations estimate finite marginal laws;
the identified intervention functional is then evaluated on those estimates.
The proof uses a finite set of coordinates relevant to the queried outcome
and a continuous-mapping step. The statement is a consistency route under
these conditions. It does not deliver a finite-bank error bound or a
generally valid confidence interval for arbitrary persistent maps, shared
calendar shocks and learned representations. Rate and central-limit
conditions must be established for the actual proposed estimator.
\citep{exec23_geometric_models}

Two apparent slips in the selected paper's extracted proof text also deserve
a bounded response. The example following its amenability assumption prints
an eventually empty intersection for a finite permutation subgroup translated
by one of its own elements. Algebra instead gives $\phi G_n=G_n$, so
$|\phi G_n\cap G_n|=|G_n|$ and the required ratio is one. Enumerating
the six permutations of three objects reproduces that identity. Separately,
the instrumental-variable proof concludes that a covariance denominator is
positive from relevance and nonzero instrument variance. Relevance warrants
nonzero, not positive, covariance. With first-stage coefficient $-2$,
instrument variance $4$ and outcome coefficient $3$, the covariance ratio
is $(-24)/(-8)=3$. Both layout and raw text extraction preserve the printed
passages. These local corrections preserve the intended algebraic routes;
they do not independently refute or fully verify the theorems. Exact
typeset-symbol inspection and the cited proof dependencies remain separate
review obligations.
\citep{exec23_geometric_models,exec23_structured_reference}

The practical evidence has its own limits. The array example explicitly
reports poor absolute uncertainty calibration despite improvement relative
to its comparators, with a mean-field approximation suggested as a cause.
In the genomic semisynthetic table, adding reverse-complement equivariance
changes the R-learner MSE from 1.637 to 2.509 in scenario A, whereas in
scenario B an outcome-model variant has MSE 0.084 compared with 0.104 for
the R-learner without that symmetry. These are author-reported experiments,
with arithmetic differences 0.872 and $-0.020$ respectively. They support
the paper's qualified discussion of interactions among model restrictions,
misspecification and estimation strategy, not a universal ranking or a bank
calibration result.
\citep{exec23_geometric_models,exec23_structured_reference}

\subsection{Hierarchical interventions need an explicit weighting target}
\label{exec23:structured-target}

Aiad and coauthors develop a complementary pipeline for hierarchical
structural causal models. Its unit-level variables represent distributions
or conditional mechanisms of subunits. For example, $Q_i^a$ describes
the treatment distribution within unit $i$. The queried contrast can be
written as a difference between expected outcomes under two interventions
on that distribution, $\mathrm{do}(Q^a=q_1)$ and
$\mathrm{do}(Q^a=q_0)$. The pipeline collapses the hierarchical graph,
augments it where required and marginalizes variables before sending the
transformed graph to pyAgrum for symbolic identification. The identification
engine does not itself inspect the raw nested table.
\citep{exec23_hierarchical_pipeline}

Even a perfectly implemented hierarchy leaves a banking target choice. Let
$n_i$ be the number of customers in branch $i$, let $\tau_i$ be a
stipulated incremental gain per customer under a specified branch policy,
and let $c$ be a uniform per-customer cost. Equal weighting of branch
effects and weighting customers equally give
\begin{equation}
 V_{\mathrm{branch}}=\frac1K\sum_{i=1}^K\tau_i-c,\qquad
 V_{\mathrm{customer}}=\frac{\sum_{i=1}^K n_i\tau_i}{\sum_{i=1}^K n_i}-c.
 \label{exec23:structured-weights}
\end{equation}
For two branches with $n_1=10$, $n_2=90$, effects USD 20 and USD $-4$,
and cost USD 2, the first target is USD 6 and the second is USD $-3.60$.
The branch-weighted quantity is an equally weighted mean of branch-specific
per-customer values; it is not total portfolio profit. Neither target is
wrong merely because the other differs. A product error occurs when one is
reported under the other's name. This is an exact authored example with
stipulated effects, not an empirical result of the hierarchical package.
\citep{exec23_structured_reference}

For an agent translating a business question, the first obligations are thus
to identify the intervention level, whether it changes a distribution or
individual treatment, the units whose outcomes are averaged, and the
sampling/weighting law. A polished diagram or a cluster-robust regression
does not answer those questions by itself. The hierarchical framework is
a candidate representation once these choices are justified; its paper
does not establish an automated business-to-model solution.

\subsection{An auditable expression tree must preserve the evaluated functional}
\label{exec23:structured-ast}

The hierarchical pipeline's abstract syntax tree records products,
marginalizations and fitted density keys. This can expose which conditional
models are required and which computations can be evaluated separately.
That is useful for a library/CLI/MCP interface: the declared estimand,
identified symbolic expression, selected numerical factors and resulting
value can remain separately inspectable. The paper reports parametric,
mixture, categorical and nonparametric fitting options and parallel
evaluation. Those descriptions are primary paper claims. The linked
repository, dependency versions, factor fallbacks and implementation have
not been newly inspected or executed in this review.
\citep{exec23_hierarchical_pipeline}

Crucially, the paper explicitly acknowledges that normalizing unstable
multiplicative factors changes the numerical functional. A constant
representation change that cancels exactly is different from a
nonconstant alteration of relative weights. To see the latter in a
transparent construction, take weights $p=(1/4,3/4)$ and effects
$\tau=(20,-4)$. Their weighted contrast is USD 2. Multiplying the two
weights by $(2,1)$ and renormalizing gives $(2/5,3/5)$, with contrast
USD $28/5=5.60$. The nominal dependency graph did not change, but the
functional did. This constructed modification is not a reconstruction of
the paper's normalization code. It shows why the exact transformation must
be declared and either proved invariant or reported as a different
functional with its own uncertainty argument.
\citep{exec23_hierarchical_pipeline,exec23_structured_reference}

The STAR application makes the distinction concrete without providing a
matched normalization experiment. Its reported mathematics row has
$2402.75-2365.95=36.80$ points. The normalized reading diagnostic has
$1788.10-1766.47=21.63$ points. These rows concern different outcomes;
their difference cannot isolate the effect of factor normalization. The
selected ExactBIC graph is described as plausible, with a substantive
orientation and added latent class context, rather than uniquely recovered.
The paper also distinguishes student-level OLS, an IV class-size coefficient
and the distributional class intervention. A larger number across these
targets is not, by itself, evidence of a more accurate causal method.
\citep{exec23_hierarchical_pipeline,exec23_structured_reference}

The stated analysis uses 322 classes with ten sampled students per class,
and identifies multi-parent factor estimation as a bottleneck. Its
conclusion leaves deeper hierarchies, uncertainty for stabilized
functionals and analysis across plausible graphs open. The reported
3.11-second sequential and 1.9-second four-process evaluations imply a
ratio of about 1.64, but this is a check of the paper's arithmetic, not a
local timing result or a bank service-level claim. These limits should
remain part of the proposed product comparison.
\citep{exec23_hierarchical_pipeline,exec23_structured_reference}

\subsection{A scoped route for the proposal}
\label{exec23:structured-route}

The two contributions suggest an orderly prototype route. First specify
the banking unit, target, information clock and intervention law. Then
justify a particular dependence restriction or hierarchical representation,
retain the assumed graph or plausible graph family, and derive the
identified functional. Preserve every numerical factor and any alteration
introduced for stability. Compare exact finite references before using
public or synthetic data to study finite behavior, uncertainty and cost.
An LLM can propose candidates, explain assumptions and help inspect traces,
but each such role needs its own evaluation and refusal rules.

For software reuse, a typed interface to a symbolic expression and its
numerical factors is worth investigating before porting an entire
implementation. That investigation must inspect the actual code, licenses,
dependencies and supported regimes. For research, the geometric restrictions
provide possible scoped assumptions; they do not replace the unresolved
joint-law and calibration arguments in the broader banking program.
Accordingly these papers strengthen the set of explicit alternatives
while leaving general discovery, arbitrary dependence and automatically
valid stabilized uncertainty as open obligations.
